% !TeX root = ../papers/nc-p3.tex
\section{十五次元成分の三十節点と二十個の二次元表現}
\label{sec:thirty-central-nodes}
\begin{theorem}\label{thm:thirty-central-nodes}
十五次元成分の一般の射影中心はちょうど三十個の通常二重点を持つ四次三次元多様体である。
二十点での単純の頂点次元は1と3、残り十点ではそれぞれ2と2である。
従って二次曲面を用いた定理\ref{thm:finite-two-dimensional-quadrics}の四辺の解は、
この一般窓につきちょうど二十個である。追加点での完備局所模型は
$R=k[[a,b,c,d]]/(ad-bc)$ と
$\operatorname{End}_R(R^{\oplus8}\oplus I^{\oplus8})$ であり、
$I$ は一般二行二列行列の余核である。
\end{theorem}
\begin{proof}
まず一般四次式の特異点でHessianの階数が4であることを、単元小行列を使う局所座標で示す。
$p=1009$ とし、節\ref{sec:unrestricted-doublet}の有理点 $t$ で、二つの隣接切断に
可逆六次小行列を取る。合計200個のSchur補成分の微分は法 $p$ で階数180を持つ。
独立な180行と180列を選ぶと、その180式は残り76変数についての陰関数を与える。
その陰関数上で残る二十式も恒等的に零である。
実際 $t$ の標数零のテンソル成分は滑らかな次元76であり、選んだ式で定まる
滑らかな同次元の局所空間の中にあるため完備局所イデアルが一致する。
$t$ のまわりで原点を移した形式陰関数は微分小行列が単元なので
$\mathbb Z_p$ 係数を持ち、自由変数の変化が $p\mathbb Z_p$ に入れば収束する。
このため有限精度の計算は、標数零の同じ成分上の実際の点の還元になる。
中心と経路空間の基底も同様である。単元小行列で取った基底についての
余分な方程式は標数零の開近傍で恒等的に成立し、積と核を与える消去も整係数で行える。

\path{computations/verify_thirty_central_nodes.py} に用いた入力と消去を全て記す。
固定した剰余ベクトルを先の微分の核へ射影し、$p$ 倍して $t$ に足した後、
選んだ180変数だけをHensel補正し全200式を法 $p^5$ で確認する。
経路長5まで順に単元階数部分を除き、四次中心の可換式の核に
元の五基底を持ち上げる。最初の四元の値を $(1,b,c,d)$ とし、
$b,c,d\ge0$, $b+c+d\le4$ の35組で第五元の左正則作用 $L_z$ を求める。
フィルター経路空間が安定した零番隅は階数64である。
四次の中心関係を $F$ とすれば特性多項式は $F^{16}$ である
（定理\ref{thm:twenty-central-nodes}のノルム恒等式）。
従って $\operatorname{tr}(L_z^j)/16$、$1\le j\le4$ とNewton公式は $F$ の係数を与える。
全次数4以下の三変数補間の行列式は単元であり、斉次化して五変数四次式を得る。

$S=k[a,b,c,d,z]$ と書き、五偏微分のイデアルを $J$、Hessianの全四次小行列の
イデアルを $H$ とする。次数9のMacaulay行列の単元ピボット数を順に示す。
各段で単元ブロックを消し、残りを $p$ で割って精度を一つ下げる。
\[
\begin{array}{c|rrrrr|r}
 &p^0&p^1&p^2&p^3&p^4&\dim S_9\\\hline
 J&637&0&48&0&0&715\\
 J+H&643&0&56&0&16&715
\end{array}
\]
これは元の行列の標数零での階数の下界を与える。各段の単元小行列は前段と独立な
非零ピボットを与えるからで、最後に精度を使い尽くした残りの階数を仮定しない。
従って $(J+H)_9=S_9$ はこの成分の稠密開集合で成立する。
特異点上ではEuler恒等式から動径がHessianの核に入るため階数はちょうど4であり、
アフィンチャートの二次主要部が非退化となる。ゆえに全特異点は通常二重点である。

次に節点 $R$ の閉点を除いた $U$ 上で任意のAzumaya類が零であることを示す。
節点の一方の小解消 $Y\to\operatorname{Spec}R$ は正則で、
$Y\setminus U=\mathbb P^1$ は余次元2である。これは二行二列行列の核直線を付ける
射影的な解消を完備局所環へ引き戻したもので、標準二チャートで正則性も分かる。
Brauer群の純性により $U$ 上の類は $Y$ に延びる。
$R$ は狭義Hensel局所環なので、固有基底変換で
$H^2_{\mathrm{et}}(Y,\mu_n)=H^2_{\mathrm{et}}(\mathbb P^1,\mu_n)=\mathbb Z/n$
となる。最後の群は核直線の射影直線の $\mathcal O(1)$ の類で生成される。
従ってKummer列から $\operatorname{Br}(Y)[n]=0$ が全 $n>0$ で成立する。
以上より完備中心の分数体 $K$ 上の一般行列環は分裂する。

その次数は16である。一般の列加群 $K^{16}$ 内に完備代数で安定な有限格子を取る
（基底の $R$ 張りを代数で生成すればよい）。局所大域次元の有限性により有限生成
射影分解を持つ。分数体へ移して列加群の重複度を数えれば、1は原始射影の
一般行列階数の整数線型結合である。従って追加点で単純が次元8の一種類しかない
可能性はない。その場合には全原始射影の一般行列階数が2だからである。
定理\ref{thm:twenty-central-nodes}の系により二単純となり、
非自由な階数1の極大Cohen--Macaulay加群による局所模型を得る。
その二行二列表示の行列式の二次主要部が非退化なので四成分はパラメータ系である。
以上で命題\ref{prop:extra-central-singular-characters}の計数の仮定が全て成立し、
$4=-56+2N$ より $N=30$、追加の単純は $2(30-20)=20$ 個となる。
\end{proof}

\begin{corollary}\label{cor:quartic-center-unramified-brauer}
一般の射影四次中心の関数体には指数4の非自明な不分岐Brauer類がある。
従ってこの三次元多様体は安定有理でない。
\end{corollary}
\begin{proof}
次数方向を除く中心の関数体で、滑らかな部分の零番隅のAzumaya類を取る。
四次中心の一次数元を可逆化すると次数方向は一つの独立な不定元になる。
その拡大先で零番隅の一般商は既知の次数4の可除環なので、拡大前も指数4である。
節点全部を吹き上げた滑らかな射影解消上で、例外因子の一般点での剰余も零となる。
実際、各節点で局所環を完備化した後の同じ吹上げでは、先の分裂により零である。
吹上げは平坦基底変換と両立し、例外因子はどちらでも同じ滑らかな接錐二次曲面で、
そのイデアルは極大イデアルの引き戻しである。よって該当する付値の分岐指数は1で
剰余体も同じであり、剰余の関手性から元の剰余は零である。
滑らかな解消上の純性から不分岐性を得る。代数閉体上の滑らかな射影多様体の
Brauer群は安定双有理不変で、射影空間では零である。
\end{proof}
